% SOURCE_AUTHORITY: sga5_fr_workpass.tex, lines 3903--3960 (LF-normalized unit).
% SOURCE_SHA256: 791F4EFFC5E02832D5D77ED03518C8156D6F07E4C8238B03545DB93D883FBB28
\subsection*{6.8}
Para ello necesitaremos recordar algunos hechos sobre la clase de cohomología ``de Hodge'' asociada a un subesquema regularmente inmerso en un esquema liso y sobre los residuos de Grothendieck ([8] III 9). Sean $f:X\to S$ un morfismo liso de dimensión relativa pura $N$ e $i:Y\hookrightarrow X$ una inmersión regular de codimensión $d$ (es decir, ((SGA~IV 16.9.2), (SGA~6 VII 1)), un subesquema cerrado definido localmente por una sucesión regular de $d$ ecuaciones o, en una terminología equivalente ((8 III 1), (SGA~6 VIII 1)), una inmersión cerrada (localmente) de intersección completa de codimensión $d$). Pongamos $g=fi$. Se asocia entonces a $Y$ un elemento
\begin{equation}
\tag{6.8.1}
\cl_{X/S}(Y)\in\Ext^d_{\cO_X}(\cO_Y,\Omega^d_{X/S}),
\end{equation}
llamado \emph{clase de cohomología asociada a $Y$}, definido de la manera siguiente. Según el teorema ``de pureza'' ([8] III 7.2), se tiene
\begin{equation}
\tag{6.8.2}
\Ext^i_{\cO_X}(\cO_Y,\Omega^d_{X/S})=
\begin{cases}
0,& i\ne d,\\
\underline{\Hom}(\bigwedge^d N_{Y/X},\Omega^d_{X/S}\otimes\cO_Y),& i=d,
\end{cases}
\end{equation}
donde $N_{Y/X}=I/I^2$ es el haz conormal ($I=\Ker\,\cO_X\to\cO_Y$). Por consiguiente,
\begin{equation}
\tag{6.8.3}
\Ext^d_{\cO_X}(\cO_Y,\Omega^d_{X/S})=H^0(Y,\underline{\Hom}(\bigwedge^dN_{Y/X},\Omega^d_{X/S}\otimes\cO_Y)).
\end{equation}

La diferencial $d_{X/S}:\OO_X\to\Omega^1_{X/S}$ induce una aplicación $\OO_Y$-lineal
\[
d_{X/S}\otimes 1:N_{Y/X}\longrightarrow \Omega^1_{X/S}\otimes\OO_Y.
\]
La $d$-ésima potencia exterior de esta aplicación es una sección global, sobre $Y$, de $\Lambda^d N_{Y/X}^{\vee}\otimes\Omega^d_{X/S}$, cuya imagen en $\Ext^d_{\OO_X}(\OO_Y,\Omega^d_{X/S})$ por el isomorfismo 6.8.3 es, por definición, la clase 6.8.1.

Supongamos que $g$ es finito y plano (por tanto, $d=N$). Sea $\omega\in H^0(Y,\Lambda^d N_{Y/X}^{\vee}\otimes\Omega^d_{X/S})$. Según 6.8.2 y 6.1, se puede considerar $\omega$ como un homomorfismo
\[
\OO_Y\longrightarrow g^!\OO_S(=\Lambda^d N_{Y/X}^{\vee}\otimes\Omega^d_{X/S}),
\]
y, por adjunción, como un homomorfismo $g_*\OO_Y\to\OO_S$. Se define el residuo (relativo a $g$) de $\omega$ como la sección de $\OO_S$ que es imagen de $1$ por este homomorfismo, y se denota por $\operatorname{Res}_{Y/S}(\omega)$ (cf. ([8] III 9)). Denotemos, por otra parte, por $i_*\omega\in H^N(X,\Omega^N_{X/S})$ la imagen de $\omega$ por la compuesta del isomorfismo 6.8.3 y de la flecha canónica
\[
\Ext^N_{\OO_X}(\OO_Y,\Omega^N_{X/S})\longrightarrow H^N(X,\Omega^N_{X/S}).
\]
Cuando $f$ es propio, de la transitividad de los morfismos de traza (o flechas de adjunción) $i_*i^!\to 1$ y $f_*f^!\to 1$ se deduce que
\begin{equation}
\tag{6.8.4}
\operatorname{Res}_{Y/S}(\omega)=\operatorname{Tr}_f(i_*\omega),
\end{equation}
donde $\operatorname{Tr}_f:H^N(X,\Omega^N_{X/S})\to H^0(S,\OO_S)$ está definido por la flecha de adjunción $f_*f^!\OO_S=Rf_*\Omega^N_{X/S}[N]\to\OO_S$ (6.1). Por otra parte, sin suponer que sea propio, se tiene la compatibilidad importante ([8] III 9 R6)
\begin{equation}
\tag{6.8.5}
\operatorname{Res}_{Y/S}(y\cl_{X/S}(Y))=\operatorname{Tr}_{Y/S}(y),
\end{equation}
para $y\in H^0(Y,\OO_Y)$, donde $\operatorname{Tr}_{Y/S}:g_*\OO_Y\to\OO_S$ designa la traza en el sentido clásico, es decir, $\operatorname{Tr}_{Y/S}(y)$ = traza del endomorfismo de $g_*\OO_Y$ definido por la multiplicación por $y$ (esta compatibilidad no se demuestra en (loc. cit.); el caso $N=1$ se trata en [Raynaud, Anneaux locaux henséliens, VII 1, Lecture Notes in Math. n\textsuperscript{o} 169, Springer Verlag, 1970]). En otros términos, $\operatorname{Tr}_{Y/S}$ corresponde, por adjunción, al homomorfismo $\OO_Y\to g^!\OO_S$ definido por la multiplicación por $\cl_{X/S}(Y)$. De ello resulta que, para $F\in\ob D(S)_{\coh}$, la flecha
\begin{equation}
\tag{6.8.6}
g^*F\longrightarrow g^!F=g^*F\otimes^{\LL}g^!\OO_S
\end{equation}
dada por el producto por $\cl_{X/S}(Y)$ corresponde, por adjunción, a la flecha
\begin{equation}
\tag{6.8.7}
g_*g^*F=F\otimes^{\LL}g_*\OO_Y\longrightarrow F
\end{equation}
definida por $\operatorname{Tr}_{Y/S}:g_*\OO_Y\to\OO_S$.

